% !TeX root = ../bnuthesis-example.tex

\chapter{文献引用}

\section{引文与参考文献的一致性}

2026 年编写规则要求引文标注与参考文献表采用 GB/T 7714-2015，在顺序编码制或著者--出版年制中选择一种，并保持全文一致。全部被引文献都应进入参考文献表，未引用的文献不应混入其中。

本例按已确认的例外，保留 \file{bnuthesis-author-year-bnu.bst} 对文献表的著者排序，并用数字上标标注正文引用。这一组合需要培养单位认可，不能作为规则规定的标准配套方式。中文文献的著者排序依赖数据库中的拼音 \texttt{key} 字段，新增条目时须补充并核对该字段。

\section{本例使用的引用方式}

本例使用 BibTeX 与 \pkg{natbib}，全文保持相同的数字上标格式。例如，期刊论文\cite{zhangkun1994}、专著\cite{zhukezhen1973} 和学位论文\cite{zhengkaiqing1987} 可以分别引用。同一位置引用多篇文献时，将键写在同一命令内\cite{zhangkun1994,zhukezhen1973}。

下面给出命令和排版结果的对应关系，各行均使用当前引用制度。
\begin{center}
  \begin{tabular}{ll}
    \toprule
    命令 & 结果 \\
    \midrule
    \verb|\cite{zhangkun1994}| & \cite{zhangkun1994} \\
    \verb|\citet{zhangkun1994}|              & \citet{zhangkun1994}              \\
    \verb|\citep{zhangkun1994}|              & \citep{zhangkun1994}              \\
    \verb|\cite[42]{zhukezhen1973}| & \cite[42]{zhukezhen1973} \\
    \verb|\cite{zhangkun1994,zhukezhen1973}| & \cite{zhangkun1994,zhukezhen1973} \\
    \bottomrule
  \end{tabular}
\end{center}

数据库中还保留了原有的其他文献条目，只有在正文中被引用的条目才输出到文献表。概率论专著\cite{peebles2001probability} 在此作为外文专著的著录示例，不表示本例进行了相应的研究。

\section{其他配套方式的代码说明}

若采用标准顺序编码制，可在导言区选用数字文献表样式，按引用顺序著录。以下仅展示代码，不在本例正文中执行：
\begin{verbatim}
\usepackage[sort]{natbib}
\bibliographystyle{bnuthesis-numeric}
\end{verbatim}

若采用标准著者--出版年制，应同时更换文献表样式和正文标注方式，并将中文文献按著者字顺和年份排列。以下也仅作代码说明：
\begin{verbatim}
\usepackage{natbib}
\bibliographystyle{bnuthesis-author-year}
\bnusetup{cite-style = author-year}
\end{verbatim}

修改参考文献方案后，应完整运行 LaTeX、BibTeX 和后续 LaTeX 编译，检查是否有未解析的引用以及正文与文献表不对应的条目。

\cite{zhangkun1994,zhukezhen1973,dupont1974bone,zhengkaiqing1987,%
  jiangxizhou1980,jianduju1994,merkt1995rotational,mellinger1996laser,%
  bixon1996dynamics,mahui1995,carlson1981two,taylor1983scanning,%
  taylor1981study,shimizu1983laser,atkinson1982experimental,%
  kusch1975perturbations,guangxi1993,huosini1989guwu,wangfuzhi1865songlun}

\cite{zhaoyaodong1998xinshidai,biaozhunhua2002tushu,chubanzhuanye2004,%
  who1970factors,peebles2001probability,baishunong1998zhiwu,%
  weinstein1974pathogenic,hanjiren1985lun,dizhi1936dizhi,%
  tushuguan1957tushuguanxue,aaas1883science,fugang2000fengsha,%
  xiaoyu2001chubanye,oclc2000about,scitor2000project%
}
